% tenkz Phase-0 gallery: benchmarks, MPO stress cases, and the
% lattice sub-language, tightly cropped.
\documentclass[varwidth=19cm, border=8pt]{standalone}
\usepackage{amsmath, amssymb}
\usepackage{tenkz}
\newcommand{\head}[1]{\par\medskip\noindent\textbf{#1}\par\smallskip}
\begin{document}

\head{B1 — periodic MPS word}
\begin{tenkz}[periodic, physical=up, bond label={$D$ at 1-2}]
  \tn[up=$i_1$]{A} & \tn[up=$i_2$]{A} & \tn[up=$i_3$]{A}
\end{tenkz}

\head{B2 — gauge equation}
% Gauge equation B^i = X A^i X^{-1} between D x D matrices: both sides
% carry the same boundary signature (2 open virtual + 1 physical); the
% protruding stubs ARE the matrix indices.
\[
  \begin{tenkz}[physical=up]
    \tn[up=$i$]{B}
  \end{tenkz}
  \;=\;
  \begin{tenkz}[physical=up]
    \tnX{X} & \tn[up=$i$]{A} & \tnX{X^{-1}}
  \end{tenkz}
\]

\head{B3 — transfer map}
% E_A = sum_i A^i (x) conj(A^i): a map on the doubled virtual space --
% four open virtual legs; west pair = input (X plugs in there), east pair
% = output.  The contracted physical leg IS the sum over i (contrast B8,
% where the sum is explicit and the leg stays open).
\[
  \mathcal{E}_A \;=\;
  \begin{tenkz}[sandwich]
    \tn{A} \\
    \tn*{A}
  \end{tenkz}
\]

\head{B4 — blocking}
% TI blocking: A^{(L)}_{(i_1...i_L)} = A^{i_1} ... A^{i_L} as D x D
% matrices -- the virtual boundary indices stay open on BOTH sides; the
% physical counts differ only by grouping (L legs vs one combined index).
\[
  \begin{tenkz}[physical=up]
    \tn[up=$i_1$]{A}\tnspan[brace below]{4}{L\text{ copies}} &
    \tn[up=$i_2$]{A} & \tndots & \tn[up=$i_L$]{A}
  \end{tenkz}
  \;=\;
  \begin{tenkz}[physical=up]
    \tn[pill, wide=2, up={$i_1\ldots i_L$}]{A^{(L)}}
  \end{tenkz}
\]

\head{B5 — fusion-tree skins}
% The pentagon-coherence layout itself belongs to tikz-cd; the trees are
% the structured atoms this package owns, in both strand skins.
The wire skin \tntree[tree style=wire]{(((a\,b)_x\,c)_y\,d)_e} and the
ribbon skin \tntree[tree style=ribbon]{(((a\,b)_x\,c)_y\,d)_e} share one
topology.

\head{B6 — PEPS window with regions (tenkzlattice)}
\begin{tenkzlattice}[rows=4, cols=4, boundary legs]
  \tnregion[slot=selected,  name=R, label=$R$]{(1-3, 1-3)}
  % R = [1,3] x [1,3];  S = R \ {(2,2)}.  S's outer boundary coincides
  % with R's, so it is drawn with inset=1 (one nesting unit tighter); the
  % ring around the removed site is S's inner boundary.
  \tnregion[slot=secondary, outline, inset=1, label=$S$, label at=south west]{R - (2,2)}
  \tnedge[distinguished]{(2,3)-(2,4)}
  \tnsite[removed]{(2,2)}
\end{tenkzlattice}

\head{B7 — zipper / pulling-through}
% Zipper (pulling through, arXiv:2203.12563): V (O-window (x) A-window)
%   = (B-window) V, with V : C^{D_O} x C^{D_A} -> C^{D_B} fusing the two
% bond spaces and B^i = sum_j V (O^{i j} (x) A^j) the fused site.  LHS:
% V's single combined leg is the open west index; the east O/A pair stays
% open.  RHS: the fused-bond B word meets V's centred west port, while the
% east face presents the separate O/A pair.
\[
  \begin{tenkz}[rows={op, ket}]
    \tnfuse[span=2, west at=center, east at={1,2}]{V} &
    \tn[mpo]{O}\tnspan[brace above]{3}{n} & \tndots & \tn[mpo]{O} \\
                       & \tn{A}                                & \tndots & \tn{A}
  \end{tenkz}
  \;=\;
  \begin{tenkz}[rows={ket:fused:nopair, wire:west none}]
    \tn[box, up=$i_1$]{B}\tnspan[brace above]{3}{n} & \tndots &
    \tn[box, up=$i_n$]{B} &
    \tnfuse[span=2, west at=center, east at={1,2}]{V}\\
    & & &
  \end{tenkz}
\]

\head{B8 — inline math atom}
% E_A(X) = sum_i A^i X (A^i)^dagger: the sum index i is explicit, so no
% contracted physical leg is drawn -- each summand is a product of three
% matrices on ONE virtual wire, open at both ends (a map, not a scalar).
Since
$\mathcal{E}_A(X)=\sum_{i}\,\tnpic[inline]{\tn[box]{A^i} & \tnX{X} & \tn*[box]{A^i}}$,
each summand is a congruence.

\head{S1 — expectation value window}
% Local window of <psi|O|psi>: physical indices contract vertically
% (ket-op, op-bra); all six virtual indices stay open -- this is a window
% of a larger contraction, not a scalar.
\[
  \begin{tenkz}[rows={ket, op, bra}]
    \tn{A} & \tn{A} & \tn{A} \\
    \tn[mpo]{O} & \tn[mpo]{O} & \tn[mpo]{O} \\
    \tn*{A} & \tn*{A} & \tn*{A}
  \end{tenkz}
\]

\head{S2 — MPO--MPO product, operator-hued wires}
\[
  \begin{tenkz}[rows={op:operator, op:operator}]
    \tn[mpo]{T_g} & \tndots & \tn[mpo]{T_g} \\
    \tn[mpo]{T_h} & \tndots & \tn[mpo]{T_h}
  \end{tenkz}
\]

\head{S3 — periodic MPO acting on a periodic MPS}
% O|psi> for a periodic MPO on a periodic MPS:
%   (O|psi>)^{i_1 i_2 i_3} = sum_j Tr[M^{i_1 j_1} M^{i_2 j_2} M^{i_3 j_3}]
%                                  Tr[A^{j_1} A^{j_2} A^{j_3}]
% Each row closes by its own virtual trace; the physical j's contract
% between the rows; the i's stay open above.
\[
  \begin{tenkz}[periodic, rows={op, ket}]
    \tn[mpo]{M} & \tn[mpo]{M} & \tn[mpo]{M} \\
    \tn{A} & \tn{A} & \tn{A}
  \end{tenkz}
\]

\head{S4 — entanglement cut}
\[
  \begin{tenkz}[physical=up]
    \tn{A} & \tn{A}\tncut{S_{[1,2]}} & \tn{A} & \tn{A}
  \end{tenkz}
\]

\head{S5 — fused wire with gauge insertions}
\[
  \begin{tenkz}[rows={wire:fused}]
    \tn[box]{B_1} & \tnX{X} & \tn[box]{B_2} & \tnX{X^{-1}} & \tn[box]{B_3}
  \end{tenkz}
\]

\head{S6 — mixed transfer with insertion}
\[
  \begin{tenkz}[sandwich]
    \tn{A} & \tnX{X} & \tn{A} \\
    \tn*{B} & \tnghost{} & \tn*{B}
  \end{tenkz}
\]

\head{T1 — standalone fusion tree in text}
A module tree \tntree{((a\,b)_x\,c)_e} sits on the math axis.

\end{document}
